\subsubsection{Anti-Landauer entropy and the index balance}

\begin{proposition}[Anti-Landauer entropy]\label{thm:anti-landauer-bits}\lean{Infomaterial.AntiLandauer.antiLandauerBits}\leanok\uses{thm:coordination-laplacian-dual}
Anti-Landauer entropy is the reversible harmonic side of the cover nerve:
$\check H^0=|C|$, the global-section count kept by settling.
\end{proposition}

\begin{proposition}[Euler characteristic of the cover nerve]\label{thm:nerve-euler-char}\lean{Infomaterial.AntiLandauer.eulerChar}\leanok\uses{thm:cover-nerve}
The Euler characteristic of the cover nerve is $\chi=|V|-|E|$ over
$\mathbb Z$.
\end{proposition}

\begin{theorem}[The index balance $\check H^0-\check H^1=\chi$]\label{thm:anti-landauer-index-balance}\lean{Infomaterial.AntiLandauer.index_balance}\leanok\uses{thm:anti-landauer-bits, thm:nerve-euler-char, thm:coordination-h1rank}
Anti-Landauer bits minus Landauer bits equals the Euler characteristic:
$\check H^0-\check H^1=\chi$. Settling preserves the index while distributing
it between harmonic content and obstruction content.
\end{theorem}

\begin{proposition}[Consensus and contradiction charges]\label{thm:consensus-contradiction-charges}\lean{Infomaterial.AntiLandauer.consensusCharge}\leanok\uses{thm:anti-landauer-bits, thm:coordination-h1rank}
The same index ledger can be read as two finite charges: consensus charge is
$\check H^0$, and contradiction charge is $\check H^1$.
\end{proposition}

\begin{theorem}[The signed charge ledger]\label{thm:net-charge-balance}\lean{Infomaterial.AntiLandauer.net_charge_balance}\leanok\uses{thm:anti-landauer-index-balance, thm:consensus-contradiction-charges}
Consensus charge minus contradiction charge is the conserved net charge
$\chi$.
\end{theorem}

\begin{theorem}[The zero-net vacuum]\label{thm:anti-landauer-vacuum}\lean{Infomaterial.AntiLandauer.vacuum_iff_charge_balance}\leanok\uses{thm:net-charge-balance}
The finite-index vacuum is exactly the state in which consensus and
contradiction charges cancel.
\end{theorem}

\begin{theorem}[Equal pair creation preserves the index]\label{thm:equal-pair-creation-preserves-net}\lean{Infomaterial.AntiLandauer.equal_pair_creation_preserves_net}\leanok\uses{thm:net-charge-balance}
Adding the same finite amount to the consensus side and to the contradiction
side preserves the net charge.
\end{theorem}

\begin{theorem}[Fixed-net cycle creation requires consensus creation]\label{thm:fixed-net-cycle-requires-consensus}\lean{Infomaterial.AntiLandauer.fixed_net_cycle_requires_consensus}\leanok\uses{thm:net-charge-balance}
At fixed net charge, increasing the contradiction side by $d$ forces the
consensus side to increase by the same $d$.
\end{theorem}

\begin{theorem}[Fixed-net consensus annihilation requires cycle annihilation]\label{thm:fixed-net-consensus-annihilation-requires-cycle-annihilation}\lean{Infomaterial.AntiLandauer.fixed_net_consensus_annihilation_requires_cycle_annihilation}\leanok\uses{thm:net-charge-balance}
At fixed net charge, decreasing the consensus side by $d$ forces the
contradiction side to decrease by the same $d$.
\end{theorem}

\begin{theorem}[Anti-descent and mirror temperature]\label{thm:anti-descent-mirror-temperature}\notready\uses{thm:fixed-net-cycle-requires-consensus, thm:positive-temperature-free-energy-flow}
The current finite theorem proves only the index-compensation law. A genuine
mirror temperature theorem must add a temporal or stochastic carrier showing
that the anti-descent builds $\check H^0$ while free-energy descent discharges
$\check H^1$.
\end{theorem}

\begin{theorem}[A reversible nerve's index is all harmonic]\label{thm:reversible-index-harmonic}\lean{Infomaterial.AntiLandauer.reversible_index_all_harmonic}\leanok\uses{thm:anti-landauer-index-balance}
When the Landauer side vanishes, the Euler characteristic is carried wholly
by the anti-Landauer harmonic side.
\end{theorem}

\begin{theorem}[The tree index balance]\label{thm:tree-index-balance}\lean{Infomaterial.AntiLandauer.tree_index_balance}\leanok\uses{thm:anti-landauer-index-balance, thm:tree-coordination-free}
For a finite tree nerve, $\check H^0=1$, $\check H^1=0$, and $\chi=1$.
This is the index balance on the connected acyclic case.
\end{theorem}
